\subsection{Fidelity of the aggregate update}\label{sec:fidelity}
The simulator draws each batch's pulses as independent Poisson counts per cell and applies them in one exact step of the
interleaved dynamics (Section~\ref{sec:model}). A physical stochastic-pulse array instead fires row and column lines per slot
and per sample; a cell is pulsed when both of its lines fire, and the samples follow one another. After the second review we
compared the two, pre-registered with pass lines (\texttt{PLAN-REVISION-2.md}, experiment A), on the operating regime of the
main table.

\paragraph{Update level (A1).} At two points of training, one batch's inputs, errors and array state were taken from surrogate
training runs and applied repeatedly from the same state by three methods: the surrogate; the explicit stream applied pulse by
pulse (Bernoulli line firings, the per-pulse device step, write noise and bounds after every pulse; the literal stream); and a
per-sample form of the same stream (seq), which replays the literal stream's line firings. Within one sample every pulse of a
cell has the same sign, so seq applies the sample's train in closed form. Its write noise is an approximation: it is added once
per sample at the pre-sample state (standard deviation $c_{2c}\sqrt{n}$ times the local step for $n$ pulses), and the bounds are
applied after each sample, where the literal stream applies noise and bounds after every pulse. The cases cover clipped and
unclipped layers: MNIST Li-ECRAM c-TTv2 (first and output layer), CIFAR MO-ECRAM AGAD (fully connected and output layer; the
worst-clipped layer of the main table) and the LM protonic-inspired c-TTv2 (embedding and output layer), \RYProbeN{} probes.
Version 3 compared seq with the literal stream as if their draws were independent, although they share the line firings, so
its sampling floor could hide a real difference (Table~\ref{tab:v3s_rev2_probe}, Appendix~\ref{app:revision}). The probes were
therefore run again with the draws kept, 16 per probe, and seq is now compared paired by draw, with bootstrap bounds over whole
draws (Tables~\ref{tab:rev4_probe_sur} and~\ref{tab:rev3_probe}; \texttt{PLAN-REVISION-3.md}); the draws are in
the deposit inventory, and Tables~\ref{tab:rev4_probe_sur}--\ref{tab:rev3_dep} can be rebuilt from them and the probe
records' metadata (Section~\ref{sec:availability}).
\begin{itemize}
\item The surrogate, whose draws are independent of the literal stream's, meets the pre-registered lines (realised gain along
      the requested update within 10\,\%, mean difference beyond sampling noise at most 0.10, median variance ratio in
      $[0.5, 2]$) on \RYProbeSurMeetsN{} of \RYProbeN{} probes, including layers that lose up to \RYProbeOkLostMax\,\% of the
      requested pulse mass to clipping; its variance ratio lies between \RYProbeSurVarMin{} and \RYProbeSurVarMax. These are
      point estimates: with 16 draws its bootstrap bounds stay within the lines on \RYProbeSurAtNinetyFiveN{} of them.
\item It fails on the two probes of the most clipped layer: \RYProbeSurFailList. There almost nothing of the requested update
      survives the clipping, and the surrogate's Poisson counts, which can exceed the slot budget that caps the stream, no longer
      match it (review 2's limiting case: Poisson(31) exceeds 31 with probability 0.4524).
\item Paired by draw, seq meets its registered lines (gain within 5\,\%, excess at most 0.05, variance ratio in
      $[0.8, 1.25]$) on \RYProbeSeqMeetsN{} of \RYProbeN{} probes. The paired floor, \RYProbeFloorRatioMin--\RYProbeFloorRatioMax{}
      times the independent one, accounts for the raw relative difference (at most \RYProbeSeqRawMax) up to an excess of at
      most \RYProbeSeqExcessMax; gains agree within \RYProbeSeqGainDevMax\,\% and variance ratios lie in
      \RYProbeSeqVarMin--\RYProbeSeqVarMax. The bootstrap bounds stay within the lines on \RYProbeSeqAtNinetyFiveN{} probes
      (\RYProbeSeqAtNinetyFiveList). On the others (\RYProbeSeqNotNinetyFiveList) 16 draws cannot exclude a mean discrepancy of
      up to \RYProbeSeqExcessPNinetyFiveMax. Seq therefore meets our registered diagnostic thresholds; that is a diagnostic,
      not a demonstration that it is equivalent to the literal stream.
\item Shared lines make the literal stream's cells dependent; the surrogate's independent draws do not reproduce this.
      Measured by per-pair correlations, which cannot cancel (Table~\ref{tab:rev3_dep}), the root-mean-square correlation among
      cells sharing a line exceeds the surrogate's sampling null by up to \RYDepExcessMax{} (mean correlation aligned with the
      expected update's signs up to \RYDepAlignedMax), by more than 0.05 on \RYDepLargeN{} of \RYProbeN{} probes
      (\RYDepLargeList). The statistic of version 3, a pooled, signed normalised covariance, stays within \RYDepVthreePooledMax{}
      in absolute value on the same draws, because positive and negative terms cancel and unequal cell variances are pooled.
      The cross-cell dependence is therefore not small everywhere; the per-cell lines above do not test it, and whether it
      matters for training is what A2 measures.
\end{itemize}
The probes draw from a generator of their own. On the pod (version 3) the MNIST and CIFAR probe runs reproduced their main-table
records exactly; on the desk (Quadro P2200) \RYProbeRunsEqualEThree{} of \RYProbeRunsN{} do, and none shares its
initialisation with its version-3 run (the MNIST run not its example order either): the simulator's generators run on
the GPU, and the CUDA random kernels assign draws to threads by the GPU's multiprocessor count. The desk's probe runs
are therefore fresh training realisations, and the probed states are not version 3's; the comparison within each
probe is unaffected, since all three updates start from the same saved state. LM training is not bitwise reproducible even on one card
type: configurations repeated exactly in the records differ by up to \RXReproLmMaxPct\,\% in bits per byte (CIFAR
\RXReproCifarMaxPct\,\%).

\paragraph{Training level (A2).} The same selections, learning rates and seeds as the main table were trained with the
per-sample stream (MNIST: every analog layer, three devices $\times$ SGD, c-TTv2, AGAD; CIFAR: the two fully connected layers,
whose rows are one per image, for Li-ECRAM and protonic-inspired c-TTv2 and MO-ECRAM AGAD; the convolutional layers keep the
surrogate), seeds 0--2, and compared seed by seed with the main table's surrogate runs; both arms are labelled at the same three
seeds (Table~\ref{tab:rev3_fid}). The pre-registered rule calls the change material if any cell's mean gap moves by more than
half the line (1\,pp) or two cells whose gaps differ by more than 1\,pp swap order.
\begin{itemize}
\item By that rule the result is \textbf{material}, through one cell: \RXFidMaterialList, whose gap moves from
      \RXFidMnistEcramMoKimTwoZeroOneNineSgdSur{} to \RXFidMnistEcramMoKimTwoZeroOneNineSgdSeq\,pp
      (difference \RXFidMnistEcramMoKimTwoZeroOneNineSgdD\,pp, 95\,\% interval [\RXFidMnistEcramMoKimTwoZeroOneNineSgdLo,
      \RXFidMnistEcramMoKimTwoZeroOneNineSgdHi]). At three seeds both of its arms are
      \RYFidMnistEcramMoKimTwoZeroOneNineSgdSurLabel{} (surrogate [\RYFidMnistEcramMoKimTwoZeroOneNineSgdSurLo,
      \RYFidMnistEcramMoKimTwoZeroOneNineSgdSurHi], stream [\RYFidMnistEcramMoKimTwoZeroOneNineSgdSeqLo,
      \RYFidMnistEcramMoKimTwoZeroOneNineSgdSeqHi]\,pp). The rerun flagged in version 3 was done as matched triples on one GPU (experiment A2-M, Table~\ref{tab:rev4_a2m};
      \texttt{PLAN-REVISION-4.md}): on the desk's \RZAtwoMGpu, each of seeds 0--4 trained the FP-SGD control, the
      surrogate and the stream from one initialisation and one example order (equal hashes, checked; the analysis refuses
      a seed without them). At these \RZAtwoMN{} matched seeds the stream's gap is \RZAtwoMSeqGap\,pp [\RZAtwoMSeqLo,
      \RZAtwoMSeqHi], \RZAtwoMSeqLabel, against \RZAtwoMSurGap\,pp [\RZAtwoMSurLo, \RZAtwoMSurHi], \RZAtwoMSurLabel,
      with the surrogate; the paired change is \RZAtwoMD\,pp [\RZAtwoMDLo, \RZAtwoMDHi] (\RZAtwoMMaterial; its interval lies within $\pm1$\,pp). The FAIL label holds under either update.
      The pod's runs of seeds 0--2 are other random realisations of the same seeds, not a rounding variant of them: the simulator's generators run on the GPU, and the CUDA random kernels assign draws to threads by the GPU's multiprocessor count, so another GPU model draws another initialisation and example order (as predicted, the hashes differ from both the pod's and the earlier desk runs').
      Per seed the two realisations' gaps differ by up to \RZRealGapDiffMax\,pp, and at the pod's three matched seeds the
      change was \RZPodNThreeD\,pp [\RZPodNThreeDLo, \RZPodNThreeDHi]: this cell's gap depends strongly on the
      realisation, and the update's effect on it is measured only by the paired comparisons within each matched set.
      Version 4's five-seed figures for this cell mixed matched and unmatched pairs and are withdrawn (Appendix~\ref{app:revision}).
\item One label changes at the same three seeds: \RYFidTransList. Its mean gap moves by only
      \RXFidCifarEcramLiTangTwoZeroOneEightCttvTwoD\,pp, but the stream's interval
      ([\RYFidCifarEcramLiTangTwoZeroOneEightCttvTwoSeqLo, \RYFidCifarEcramLiTangTwoZeroOneEightCttvTwoSeqHi]\,pp) is wider and
      includes the line. This is not a failure, but the label is not preserved, and the difference interval
      ([\RXFidCifarEcramLiTangTwoZeroOneEightCttvTwoLo, \RXFidCifarEcramLiTangTwoZeroOneEightCttvTwoHi]\,pp) does not bound
      the change within $\pm1$\,pp.
\item \RYFidRevN{} pairs of cells reverse their mean order (\RYFidRevList); none was separated by more than 1\,pp under the
      surrogate, so there is no material rank reversal under the registered rule, but the ranking as a whole is not preserved.
      No Tiki-Taka-family cell's mean moves by more than \RXFidTtMaxAbs\,pp (MNIST at most \RXFidMnistMaxAbs\,pp including the
      SGD cell; CIFAR at most \RXFidCifarMaxAbs\,pp, with intervals of several pp at three seeds).
\end{itemize}
Within these bounds (MNIST on every layer, CIFAR on its fully connected layers, the LM at update level only, three seeds), the
registered result is one material mean change, in a failing SGD cell whose FAIL label holds at five matched seeds on one GPU under either update; no
material rank reversal; no Tiki-Taka-family mean shift above \RXFidTtMaxAbs\,pp; and one label that loses PASS. This is not a
demonstration that the surrogate and the stream train equivalently, and the per-sample stream is itself an approximation of the
literal one. The surrogate is not validated where clipping removes nearly all of the update (the CIFAR output layer above), for
the convolutional layers, or for the LM's training trajectory, and the main results remain comparisons within this
aggregate-update simulator.

% generated by revision4/analyze_rev4.py (do not edit)
\begin{table}[t]\centering\footnotesize\setlength{\tabcolsep}{3pt}\caption{Update fidelity of the aggregate surrogate against the literal pulse stream (experiment A1-P; desk rerun of A1 with the draws kept, 16 draws per probe; Li = Li-ECRAM, MO = MO-ECRAM, Prot. = protonic-inspired). Lost = requested pulse mass removed by the probability clipping (\%). The surrogate's draws are independent of the stream's, so the independent sampling floor applies: gain difference (\%, bootstrap 95 \% interval over draws), excess relative mean difference beyond the floor (bootstrap 95th percentile), median variance ratio, and whether the point estimates meet the registered lines (10 \%, 0.10, [0.5, 2]). ``(95)'': the bootstrap 95th percentile of the excess and the bootstrap 95 \% gain interval are also within their lines; it does not cover the variance ratio (a point estimate) and is not a joint 95 \% guarantee for every line.}\label{tab:rev4_probe_sur}
\fitbox{\begin{tabular}{lrrrlllc}\toprule Case & L. & Upd. & Lost & $\Delta g$ \% & excess & var. & meets\\\midrule
CIFAR MO AGAD & 3 & 600 & 36.2 & 1.4 [1.2, 1.6] & 0.06 (0.17) & 1.14 & yes\\
CIFAR MO AGAD & 4 & 600 & 95.7 & 1.5 [0.5, 2.6] & 0.15 (0.16) & 1.26 & no\\
CIFAR MO AGAD & 3 & 2000 & 26.2 & 0.5 [-0.3, 1.3] & 0.03 (0.31) & 1.07 & yes\\
CIFAR MO AGAD & 4 & 2000 & 97.3 & -4.2 [-7.8, -0.2] & 0.30 (0.32) & 1.27 & no\\
LM Prot. c-TTv2 & 0 & 750 & 6.1 & 1.3 [0.8, 1.8] & 0.03 (0.14) & 1.01 & yes\\
LM Prot. c-TTv2 & 9 & 750 & 0.0 & -0.5 [-2.4, 1.8] & 0.03 (0.32) & 0.86 & yes\\
LM Prot. c-TTv2 & 0 & 2250 & 7.0 & -0.6 [-1.0, -0.1] & 0.04 (0.13) & 1.05 & yes\\
LM Prot. c-TTv2 & 9 & 2250 & 0.0 & 0.3 [-1.8, 2.3] & 0.04 (0.32) & 0.78 & yes\\
MNIST Li c-TTv2 & 0 & 1000 & 16.0 & -0.4 [-0.8, 0.1] & 0.00 (0.14) & 1.06 & yes\\
MNIST Li c-TTv2 & 2 & 1000 & 4.6 & -1.1 [-2.6, 0.6] & 0.02 (0.10) & 1.19 & yes\\
MNIST Li c-TTv2 & 0 & 3500 & 37.5 & 0.6 [0.0, 1.1] & 0.00 (0.15) & 1.11 & yes\\
MNIST Li c-TTv2 & 2 & 3500 & 5.7 & 0.7 [-0.5, 2.1] & 0.02 (0.09) & 1.41 & yes (95)\\
\bottomrule\end{tabular}}\end{table}

% generated by revision4/analyze_rev4.py (do not edit)
\begin{table}[t]\centering\footnotesize\setlength{\tabcolsep}{3pt}\caption{Update fidelity of the per-sample stream (seq) against the literal pulse stream, paired by draw (shared line firings; experiment A1-P, the probes of Table~\ref{tab:rev4_probe_sur}): gain difference (\%, bootstrap 95 \% interval), raw relative mean difference, independent / paired sampling floor, excess beyond the PAIRED floor (bootstrap 95th percentile), median variance ratio, and whether the point estimates meet the registered lines (5 \%, 0.05, [0.8, 1.25]). ``(95)'': the bootstrap 95th percentile of the excess and the bootstrap 95 \% gain interval are also within their lines; it does not cover the variance ratio (a point estimate) and is not a joint 95 \% guarantee for every line.}\label{tab:rev3_probe}
\fitbox{\begin{tabular}{lrrlllllc}\toprule Case & L. & Upd. & $\Delta g$ \% & raw & floor ind./pair. & excess & var. & meets\\\midrule
CIFAR MO AGAD & 3 & 600 & 0.1 [-0.1, 0.2] & 0.114 & 0.121 / 0.114 & 0.011 (0.144) & 1.00 & yes\\
CIFAR MO AGAD & 4 & 600 & -0.9 [-1.5, -0.3] & 0.050 & 0.050 / 0.050 & 0.005 (0.064) & 1.03 & yes\\
CIFAR MO AGAD & 3 & 2000 & -0.2 [-0.6, 0.1] & 0.224 & 0.241 / 0.224 & 0.015 (0.286) & 1.00 & yes\\
CIFAR MO AGAD & 4 & 2000 & -1.4 [-5.5, 3.5] & 0.072 & 0.074 / 0.073 & 0.000 (0.094) & 1.00 & yes\\
LM Prot. c-TTv2 & 0 & 750 & -0.0 [-0.2, 0.2] & 0.032 & 0.106 / 0.032 & 0.000 (0.041) & 1.00 & yes (95)\\
LM Prot. c-TTv2 & 9 & 750 & 0.0 [-0.3, 0.3] & 0.075 & 0.258 / 0.075 & 0.000 (0.095) & 1.00 & yes\\
LM Prot. c-TTv2 & 0 & 2250 & 0.1 [-0.0, 0.3] & 0.029 & 0.093 / 0.028 & 0.004 (0.037) & 1.00 & yes (95)\\
LM Prot. c-TTv2 & 9 & 2250 & 0.2 [-0.2, 0.5] & 0.073 & 0.249 / 0.073 & 0.007 (0.093) & 1.00 & yes\\
MNIST Li c-TTv2 & 0 & 1000 & -0.0 [-0.0, 0.0] & 0.036 & 0.109 / 0.036 & 0.001 (0.045) & 1.00 & yes (95)\\
MNIST Li c-TTv2 & 2 & 1000 & 0.1 [-0.1, 0.2] & 0.027 & 0.074 / 0.027 & 0.000 (0.035) & 1.02 & yes (95)\\
MNIST Li c-TTv2 & 0 & 3500 & -0.0 [-0.1, 0.0] & 0.036 & 0.113 / 0.036 & 0.003 (0.046) & 1.00 & yes (95)\\
MNIST Li c-TTv2 & 2 & 3500 & 0.1 [-0.2, 0.3] & 0.024 & 0.063 / 0.025 & 0.000 (0.031) & 0.99 & yes (95)\\
\bottomrule\end{tabular}}\end{table}

% generated by revision4/analyze_rev4.py (do not edit)
\begin{table*}[t]\centering\footnotesize\setlength{\tabcolsep}{3pt}\caption{Cross-cell dependence of the literal stream's draws (experiment A1-P), among cells sharing a row / a column, from per-pair Pearson correlations over all pairs (they cannot cancel): root-mean-square correlation of the literal stream and of the surrogate (whose cells are independent, so its value is the sampling null at the same number of draws), the excess RMS beyond that null, taken in quadrature ($\sqrt{\max(0,\,\mathrm{RMS}_{\mathrm{slot}}^2-\mathrm{RMS}_{\mathrm{null}}^2)}$), the mean correlation aligned with the signs of the expected update, the largest per-line RMS, and the version-3 statistic (a pooled, signed normalised covariance, which can cancel).}\label{tab:rev3_dep}
\fitbox{\begin{tabular}{lrrcccccc}\toprule Case & L. & Upd. & RMS slot & RMS sur.\ (null) & excess RMS & aligned mean & max line RMS & v3 pooled\\\midrule
CIFAR MO AGAD & 3 & 600 & 0.258 / 0.261 & 0.258 / 0.258 & 0.009 / 0.041 & 0.005 / 0.010 & 0.262 / 1.000 & 0.000 / -0.000\\
CIFAR MO AGAD & 4 & 600 & 0.256 / 0.261 & 0.259 / 0.260 & 0.000 / 0.026 & 0.007 / -0.004 & 0.259 / 0.645 & 0.000 / -0.008\\
CIFAR MO AGAD & 3 & 2000 & 0.260 / 0.260 & 0.258 / 0.258 & 0.032 / 0.026 & 0.007 / 0.004 & 0.283 / 0.399 & -0.000 / -0.000\\
CIFAR MO AGAD & 4 & 2000 & 0.259 / 0.260 & 0.259 / 0.255 & 0.015 / 0.048 & 0.006 / -0.003 & 0.265 / 0.315 & 0.000 / 0.001\\
LM Prot. c-TTv2 & 0 & 750 & 0.259 / 0.258 & 0.258 / 0.258 & 0.015 / 0.008 & -0.000 / 0.000 & 0.266 / 0.264 & 0.000 / 0.000\\
LM Prot. c-TTv2 & 9 & 750 & 0.299 / 0.258 & 0.258 / 0.258 & 0.152 / 0.007 & 0.026 / -0.000 & 1.000 / 0.264 & -0.000 / -0.001\\
LM Prot. c-TTv2 & 0 & 2250 & 0.258 / 0.258 & 0.258 / 0.258 & 0.000 / 0.000 & -0.000 / 0.000 & 0.265 / 0.262 & 0.000 / -0.000\\
LM Prot. c-TTv2 & 9 & 2250 & 0.315 / 0.258 & 0.258 / 0.258 & 0.180 / 0.003 & 0.025 / -0.000 & 1.000 / 0.266 & 0.001 / -0.001\\
MNIST Li c-TTv2 & 0 & 1000 & 0.335 / 0.262 & 0.258 / 0.258 & 0.213 / 0.044 & 0.054 / 0.010 & 0.454 / 0.589 & -0.000 / -0.000\\
MNIST Li c-TTv2 & 2 & 1000 & 0.399 / 0.262 & 0.258 / 0.260 & 0.305 / 0.031 & 0.103 / 0.014 & 0.494 / 0.348 & -0.000 / -0.020\\
MNIST Li c-TTv2 & 0 & 3500 & 0.385 / 0.259 & 0.258 / 0.258 & 0.286 / 0.025 & 0.100 / 0.004 & 0.478 / 0.322 & -0.001 / -0.000\\
MNIST Li c-TTv2 & 2 & 3500 & 0.431 / 0.264 & 0.258 / 0.258 & 0.345 / 0.057 & 0.208 / 0.003 & 0.522 / 0.321 & 0.006 / -0.003\\
\bottomrule\end{tabular}}\end{table*}

% generated by revision4/analyze_rev4.py (do not edit)
\begin{table}[t]\centering\footnotesize\setlength{\tabcolsep}{3pt}\caption{Experiment A2-M: MNIST MO-ECRAM SGD as matched triples on one GPU (Quadro P2200): per seed, the FP-SGD control's, the surrogate's and the per-sample stream's test accuracy (\%) and the two gaps to the control (pp); every seed's three runs share their initialisation and example order (equal hashes, checked). Right: the same seeds' gaps from the pod's matched triples (RTX A4000), a different random realisation of each seed.}\label{tab:rev4_a2m}
\fitbox{\begin{tabular}{rrrrllll}\toprule & \multicolumn{5}{c}{P2200} & \multicolumn{2}{c}{pod (A4000)}\\\cmidrule(lr){2-6}\cmidrule(lr){7-8} Seed & FP-SGD & sur. & seq & gap sur. & gap seq & gap sur. & gap seq\\\midrule
0 & 97.90 & 89.83 & 90.02 & 8.07 & 7.88 & 16.84 & 15.69\\
1 & 97.60 & 87.86 & 87.30 & 9.74 & 10.30 & 7.59 & 7.39\\
2 & 97.68 & 89.79 & 90.26 & 7.89 & 7.42 & 17.86 & 14.95\\
3 & 97.74 & 87.20 & 87.22 & 10.54 & 10.52 & -- & --\\
4 & 97.75 & 90.44 & 90.20 & 7.31 & 7.55 & -- & --\\
\midrule mean [95\,\%] & & & & 8.71 [7.02, 10.40] & 8.73 [6.82, 10.65] & & \\
\bottomrule\end{tabular}}\end{table}

% generated by revision3/analyze_rev3.py (do not edit)
\begin{table*}[t]\centering\footnotesize\caption{Training fidelity (experiment A2), seeds 0--2 on both arms: mean gap to FP-SGD (pp) with its 95 \% interval and label against the 2\,pp line (P pass, U unresolved, F fail), with the aggregate surrogate (E3's records) and with the per-sample pulse stream (MNIST: every analog layer; CIFAR: the two fully connected layers); the paired difference with its 95 \% interval; material by the registered rule (|mean| $>$ 1\,pp). The labels are at $n=3$; the main table's labels are at $n=5$.}\label{tab:rev3_fid}
\fitbox{\begin{tabular}{lllll c}\toprule Task & Cell & surrogate & stream & difference & material\\\midrule
CIFAR & Li-ECRAM c-TTv2 & -1.14 [-3.89, 1.62] P & -0.62 [-6.44, 5.20] U & 0.52 [-3.26, 4.29] & no\\
CIFAR & MO-ECRAM AGAD & -1.40 [-3.22, 0.41] P & -1.82 [-3.76, 0.13] P & -0.41 [-1.99, 1.16] & no\\
CIFAR & Protonic-inspired c-TTv2 & -0.85 [-3.33, 1.64] P & -1.46 [-3.94, 1.02] P & -0.61 [-2.86, 1.64] & no\\
MNIST & Li-ECRAM AGAD & 0.65 [-0.26, 1.56] P & 0.49 [0.08, 0.90] P & -0.16 [-0.69, 0.36] & no\\
MNIST & Li-ECRAM c-TTv2 & 0.68 [-0.77, 2.12] U & 0.74 [-1.70, 3.18] U & 0.07 [-0.95, 1.09] & no\\
MNIST & Li-ECRAM SGD & 6.99 [0.52, 13.46] U & 6.90 [1.03, 12.76] U & -0.09 [-0.73, 0.54] & no\\
MNIST & MO-ECRAM AGAD & 0.77 [0.48, 1.06] P & 0.89 [0.70, 1.07] P & 0.11 [-0.07, 0.30] & no\\
MNIST & MO-ECRAM c-TTv2 & 1.08 [0.44, 1.72] P & 1.05 [0.43, 1.66] P & -0.04 [-0.23, 0.16] & no\\
MNIST & MO-ECRAM SGD & 14.10 [0.04, 28.15] U & 12.68 [1.27, 24.09] U & -1.42 [-4.84, 2.00] & yes\\
MNIST & Protonic-inspired AGAD & 0.79 [0.15, 1.43] P & 0.77 [0.17, 1.37] P & -0.02 [-0.39, 0.35] & no\\
MNIST & Protonic-inspired c-TTv2 & 0.29 [-0.30, 0.88] P & 0.21 [-0.71, 1.13] P & -0.08 [-0.44, 0.27] & no\\
MNIST & Protonic-inspired SGD & 4.58 [0.15, 9.00] U & 4.69 [0.01, 9.38] U & 0.12 [-0.33, 0.56] & no\\
\bottomrule\end{tabular}}\end{table*}

